\documentclass[11pt]{article}
\usepackage{fontspec}
\usepackage[french]{babel}
\usepackage{amsmath}
\usepackage{amssymb}
\usepackage{amsthm}
\usepackage{mathrsfs}
\usepackage{tikz-cd}
\usepackage{enumitem}

\newcommand{\printedpage}[1]{\marginpar{\footnotesize p.~#1}}
\newcommand{\msp}[1]{\marginpar{\footnotesize [#1]}}
\newcommand{\circled}[1]{\tikz[baseline=(c.base)]\node[draw,circle,inner sep=1pt](c){#1};}
\newcommand{\ovalbox}[1]{\tikz[baseline=(c.base)]\node[draw,rounded corners=6pt,inner sep=2pt](c){#1};}
\newcommand{\Ob}{\operatorname{Ob}}
\newcommand{\card}{\operatorname{card}}

\begin{document}

\printedpage{119}
% head: A. Grothendieck : Dérivateurs, Ch. XVIII, Édition des Universités de Montpellier II et Paris VII
satisfaisant aux contraintes $(**)$, $(***)$, $(****)$. On va faire la construction en trois étapes.

\circled{1} Départ :
\[
i_0 \;=\; i\ .
\]
Comme $J'$ est cofinal dans $J$, existe $j_0\in\Ob J'$ et
\[
\varphi(i_0)\xrightarrow{u_0} j_0\ .
\]

\msp{page 316}

\circled{2} Supposons que les $i_\beta$, $j_\beta$ soient construits pour $\beta<\alpha$, ainsi que les $\tau_{\beta',\beta}:i_\beta\longrightarrow i_{\beta'}$, $\tau'_{\beta',\beta}:j_\beta\longrightarrow j_{\beta'}$ pour $\beta,\beta'<\alpha$, enfin les $u_\beta:\varphi(i_\beta)\longrightarrow j_\beta$, $v_\beta:j_\beta\longrightarrow\varphi(i_{\beta+1})$ pour $\beta<\alpha$ resp. pour $\beta+1<\alpha$, de façon à satisfaire aux contraintes $(**)$, $(****)$, i.e. commutativités dans
\[
\left\{
\begin{array}{l}
\left.
\begin{array}{ll}
\circled{1} &
\begin{tikzcd}[ampersand replacement=\&]
\varphi(i_\beta) \arrow[r, "u_\beta"] \arrow[rr, bend right=15, "\varphi(\tau_{\beta+1,\beta})"'] \& j_\beta \arrow[r, "v_\beta"] \& \varphi(i_{\beta+1})
\end{tikzcd}
\\[5ex]
\circled{2} &
\begin{tikzcd}[ampersand replacement=\&]
j_\beta \arrow[r, "v_\beta"] \arrow[rr, dotted, bend right=15, "\tau'_{\beta+1,\beta}"'] \& \varphi(i_{\beta+1}) \arrow[r, "u_{\beta+1}"] \& j_{\beta+1}
\end{tikzcd}
\end{array}
\right\}\quad \textit{si}\ \beta+1<\alpha
\\[8ex]
\begin{array}{ll}
\circled{3} &
\begin{tikzcd}[ampersand replacement=\&]
\varphi(i_\beta) \arrow[r, "u_\beta"] \arrow[d, "\varphi(\tau_{\beta',\beta})"'] \& j_\beta \arrow[d, "\tau'_{\beta',\beta}"] \\
\varphi(i_{\beta'}) \arrow[r, "u_{\beta'}"] \& j_{\beta'}
\end{tikzcd}
\qquad\qquad \text{si } \beta<\beta'<\alpha
\end{array}
\\[10ex]
\begin{array}{ll}
\circled{4} &
\begin{tikzcd}[ampersand replacement=\&]
j_\beta \arrow[r, "v_\beta"] \arrow[d, "\tau'_{\beta',\beta}"'] \& \varphi(i_{\beta+1}) \arrow[d, "\varphi(\tau_{\beta'+1,\beta+1})"] \\
j_{\beta'} \arrow[r, "v_{\beta'}"] \& \varphi(i_{\beta'+1})
\end{tikzcd}
\qquad\qquad \text{si } \beta<\beta' \text{ et } \beta'+1<\alpha\ ,
\end{array}
\end{array}
\right.
\]

\msp{page 317}

il faut construire
% foot: Version 17 décembre 2025

\printedpage{120}
% head: A. Grothendieck : Dérivateurs, Ch. XVIII, Édition des Universités de Montpellier II et Paris VII
\begin{itemize}[leftmargin=*]
\item[1)] $i_\alpha\in\Ob I$,\quad et un homomorphisme $(i_\beta)_{\beta\in I_\alpha}\longrightarrow i_\alpha$, i.e. des $\tau_{\alpha,\beta}$ : $i_\beta\longrightarrow i_\alpha$ $(\beta\in I_\alpha)$, compatibles aux translations $\tau_{\beta',\beta}$ pour $\beta<\beta'<\alpha$.
\item[2)] $j_\alpha\in\Ob J'$,\quad et un homomorphisme $(j_\beta)_{\beta\in I_\alpha}\longrightarrow j_\alpha$, i.e. des $\tau'_{\alpha,\beta}$ : $j_\beta\longrightarrow j_\alpha$ $(\beta\in I_\alpha)$, compatibles aux translations $\tau'_{\beta',\beta}$ pour $\beta<\beta'<\alpha$.
\item[3)] Une flèche
\[
u_\alpha:\varphi(i_\alpha)\longrightarrow j_\alpha\ .
\]
\item[4)] Si $\alpha$ n'est pas ordinal limite, i.e. $\alpha=\alpha_0+1$, une flèche
\[
v_{\alpha_0}:j_{\alpha_0}\longrightarrow\varphi(i_{\alpha_0+1})=\varphi(i_\alpha)\ ;
\]
\end{itemize}

et ceci de façon à satisfaire aux contraintes \circled{1}\,\circled{2}\,\circled{3}\,\circled{4} quand on y suppose $\beta+1\leqslant\alpha$ (au lieu de $\beta'<\alpha$) dans \circled{1} et \circled{2}, et $\beta'\leqslant\alpha$ (au lieu de $\beta'<\alpha$) dans \circled{3}, et $\beta'+1\leqslant\alpha$ (au lieu de $\beta'+1<\alpha$) dans \circled{4}. Donc pour les contraintes $1^\circ$, $2^\circ$, $4^\circ$, il n'y a en tenir compte \textit{que} si $\alpha$ n'est pas ordinal limite, donc $\alpha=\alpha_0+1$, et en y faisant % [sic] il n'y a en tenir -> il n'y a à en tenir, the preposition « à » appears to be missing
\[
\beta' \;=\; \alpha_0\ .
\]
Quant à la contrainte $3^\circ$, il faut en
\msp{page 318}
tenir compte pour $\beta'=\alpha$. Donc au total, il faut exiger
\[
(3_\alpha)(\beta)\qquad\qquad
\begin{tikzcd}
\varphi(i_\beta) \arrow[r, "u_\beta"] \arrow[d, "\varphi(\circled{$\tau_{\alpha,\beta}$})"'] & j_\beta \arrow[d, "\tau'_{\alpha,\beta}" {draw, circle, inner sep=1pt}] \\
\varphi(i_\alpha) \arrow[r, "u_\alpha" {draw, circle, inner sep=1pt}] & j_\beta % [sic] j_\beta -> j_\alpha, the target of \tau'_{\alpha,\beta} and u_\alpha is j_\alpha
\end{tikzcd}
\qquad \text{commutatif pour } \beta<\alpha\ (\footnotemark[140]),
\]
\footnotetext[140]{\textbf{N.B.} Il suffit de l'exiger pour $\beta=\alpha_0$, dans le cas $\alpha=\alpha_0+1$, cf. ci-dessous, p. 323.}
(par restriction sur $\alpha$), et si $\alpha=\alpha_0+1$ est ordinal non limite, il faut exiger de plus
\[
(4_\alpha)(\beta)\qquad\qquad
\begin{tikzcd}
j_\beta \arrow[r, "v_\beta"] \arrow[d, "\tau'_{\alpha_0,\beta}"'] & \varphi(i_{\beta+1}) \arrow[d, "\varphi(\circled{$\tau_{\alpha,\beta+1}$})"] \\
j_{\alpha_0} \arrow[r, "v_{\alpha_0}" {draw, circle, inner sep=1pt}] & \varphi(i_{\alpha_0+1})=\varphi(i_\alpha)
\end{tikzcd}
\qquad \beta<\alpha_0
\]
\[
(1_\alpha)\qquad\qquad
\begin{tikzcd}
\varphi(i_{\alpha_0}) \arrow[r, "u_{\alpha_0}"] \arrow[rr, bend right=15, "\varphi(\circled{$\tau_{\alpha,\alpha_0}$})"'] & j_{\alpha_0} \arrow[r, "v_{\alpha_0}" {draw, circle, inner sep=1pt}] & \varphi(i_{\alpha_0+1})=\varphi(i_\alpha)
\end{tikzcd}
\]
% foot: Version 17 décembre 2025

\printedpage{121}
% head: A. Grothendieck : Dérivateurs, Ch. XVIII, Édition des Universités de Montpellier II et Paris VII
\[
(2_\alpha)\qquad\qquad
\begin{tikzcd}
j_{\alpha_0} \arrow[r, "v_{\alpha_0}" {draw, circle, inner sep=1pt}] \arrow[rr, bend right=15, "\tau'_{\alpha,\alpha_0}"' {draw, circle, inner sep=1pt}] & \varphi(i_{\alpha_0+1})=\varphi(i_\alpha) \arrow[r, "u_\alpha" {draw, circle, inner sep=1pt}] & j_\alpha\ .
\end{tikzcd}
\]

Dans ces compatibilités, j'ai cerclé les flèches $\tau_{\alpha,\beta}$, $\tau'_{\alpha,\beta}$ $(\beta<\alpha)$, $\Bigl(\ \tau_{\alpha,\beta+1}\ (\beta+1<\alpha)\ \Bigr)$, $v_{\alpha_0}$, $\Bigl(\ \tau_{\alpha,\alpha_0}\ \Bigr)$, $u_\alpha$, $\Bigl(\ \tau'_{\alpha,\alpha_0}\ \Bigr)$, qui ont été introduites dans les donnés 1) 2) 3) 4), page 317. % [sic] donnés -> données, the noun « données » is meant

\msp{page 319}

On va distinguer donc les deux cas.

\textit{2 A)} $\alpha=\alpha_0+1$.

Pour construire $i_\alpha$ et les $\tau_{\alpha,\beta}:i_\beta\longrightarrow i_\alpha$, il suffit de construire $i_\alpha$ et
\[
\circled{$\tau_{\alpha,\alpha_0}$}\ :\ i_{\alpha_0}\ \longrightarrow\ i_\alpha=i_{\alpha_0+1}\ .
\]
De même pour $j_\alpha$ et les $\tau'_{\alpha,\beta}$, il suffit de construire
\[
\circled{$\tau'_{\alpha,\alpha_0}$}\ :\ j_{\alpha_0}\ \longrightarrow\ j_\alpha\ .
\]
Il faut de plus se donner
\[
j_{\alpha_0}\ \xrightarrow{\circled{$v_{\alpha_0}$}}\ \varphi(\underbrace{\,i_\alpha\,}_{=i_{\alpha_0+1}})\ \xrightarrow{\circled{$u_\alpha$}} j_\alpha\ ,
\]
donc en tout deux objets $i_\alpha$, $j_\alpha$ (dans $I$ et dans $J'$ respectivement) et quatre flèches $\tau_{\alpha,\alpha_0}$, $\tau'_{\alpha,\alpha_0}$, $v_{\alpha_0}$, $u_\alpha$. Il faut quatre compatibilités (page 318), où les deux dernières $3_\alpha(\beta)$, $4_\alpha(\beta)$ dépendent d'un paramètre $\beta<\alpha$, i.e. $\beta\leqslant\alpha_0$, ou de $\beta<\alpha_0$.

\msp{page 320 biffé}

\msp{page 321}

On va d'abord s'occuper de deux premières compatibilités $1_\alpha$, $2_\alpha$, en termes de la donnée de
\[
i_{\alpha_0}\ \in\ \Ob I\ ,\qquad j_{\alpha_0}\ \in\ \Ob J'\ ,\qquad u_{\alpha_0}\ :\ \varphi(i_{\alpha_0})\ \longrightarrow\ j_{\alpha_0}\ .
\]
Les deux compatibilités voulues pour les deux objets $i_\alpha$, $j_\alpha$ et les quatre flèches $\tau_{\alpha,\alpha_0}$, $\tau'_{\alpha,\alpha_0}$, $v_{\alpha_0}$, $u_\alpha$ sont indiquées dans le diagramme
\[
(*)\qquad\qquad
\begin{tikzcd}
\varphi(i_{\alpha_0}) \arrow[r, "u_{\alpha_0}"] \arrow[rr, bend right=15, "\varphi(\tau_{\alpha,\alpha_0})"'] & j_{\alpha_0} \arrow[r, "v_{\alpha_0}\ ?"] \arrow[rr, bend left=20, "\tau'_{\alpha,\alpha_0}"] & \varphi(\underbrace{\,i_\alpha\,}_{?}) \arrow[r, "u_\alpha\ ?"] & j_\alpha\ .
\end{tikzcd}
\]

On note que comme $\varphi$ est cofinal, $\exists\ i_\alpha=i\in\Ob I$ et $j_{\alpha_0}\xrightarrow{v_{\alpha_0}}\varphi(i)$. Mais on veut que le composé
\[
w\ :\ \varphi(i_{\alpha_0})\ \xrightarrow{u_{\alpha_0}}\ j_{\alpha_0}\ \xrightarrow{v_{\alpha_0}}\ \varphi(i)
\]
soit de la forme $\varphi(\tau)$, où $\tau=\tau_{\alpha,\alpha_0}:i_{\alpha_0}\longrightarrow i_\alpha$ est une flèche dans $I$. Comme $\varphi$ est un foncteur cofinal de catégories filtrantes, on sait
% foot: Version 17 décembre 2025
\printedpage{122}
% head: A. Grothendieck : Dérivateurs, Ch. XVIII, Édition des Universités de Montpellier II et Paris VII
\msp{page 322}
qu'on peut trouver $i\xrightarrow{\rho}i'$ dans $I$, tel que la flèche composée $\varphi(i_{\alpha_0})\xrightarrow{w}\varphi(i)\xrightarrow{\varphi(\rho)}\varphi(i')$ soit de la forme $\varphi(\tau)$. Donc quitte à remplacer $i$ par $i'$, et $v_{\alpha_0}$ par le composé $v'_{\alpha_0}$ : $j_{\alpha_0}\xrightarrow{v_{\alpha_0}}\varphi(i)\xrightarrow{\varphi(\rho)}\varphi(i')$, on peut supposer que $v_{\alpha_0}u_{\alpha_0}$ est de la forme $\varphi(\tau)$. Cela nous donne donc en principe l'objet $i=i_\alpha$, et les flèches $v_{\alpha_0}:j_{\alpha_0}\longrightarrow\varphi(i)$ et $\tau_{\alpha,\alpha_0}:i_{\alpha_0}\longrightarrow i_\alpha$, de façon à satisfaire à la première compatibilité dans $(*)$. Il reste à construire $j_\alpha$, $\tau'_{\alpha,\alpha_0}$, $u_\alpha$ de façon à satisfaire à la deuxième compatibilité $(*)$. Or $J'$ étant cofinal dans $J$, il existe $j_\alpha$ dans $J'$ et une flèche $u_\alpha:\varphi(i_\alpha)\longrightarrow j_\alpha$. D'autre part, comme $J'\longrightarrow J$ est pleinement fidèle, la flèche composée $j_{\alpha_0}\xrightarrow{v_{\alpha_0}}\varphi(i_\alpha)\xrightarrow{u_\alpha}j_\alpha$ dans $J$, entre objets de $J'$, provient d'une unique flèche $\tau'_{\alpha,\alpha_0}:j_{\alpha_0}\longrightarrow j_\alpha$ dans $J'$.

\msp{page 323}

Ainsi, on a construit les deux objets et les quatre flèches voulues, de façon à satisfaire aux contraintes $1_\alpha$, $2_\alpha$ de la page 318. Les contraintes $3_\alpha(\beta)$, $4_\alpha(\beta)$ sont-elles alors satisfaites automatiquement ?
\[
\begin{tikzcd}
\varphi(i_\beta) \arrow[r, "u_\beta"] \arrow[d] & j_\beta \arrow[d] \\
\varphi(i_{\alpha_0}) \arrow[r, "u_{\alpha_0}"] \arrow[d, "\varphi(\tau_{\alpha,\alpha_0})"'] & j_{\alpha_0} \arrow[d, "\tau'_{\alpha,\alpha_0}"] \\
\varphi(i_\alpha) \arrow[r, "u_\alpha"] & j_\alpha
\end{tikzcd}
\qquad\qquad \beta<\alpha, \text{ i.e. } \beta\leqslant\alpha_0\ .
\]
Le carré supérieur est commutatif par hypothèse de récurrence, pour que le carré total (composé) soit commutatif, il suffit que le carré inférieur le soit. On a donc une double-flèche $(\tau'_{\alpha,\alpha_0}u_{\alpha_0},u_\alpha\varphi(\tau_{\alpha,\alpha_0})):\varphi(i_\alpha)\rightrightarrows j_\alpha$, elle n'est peut-être pas triviale, mais on sait, $J'$ étant filtrante, qu'on peut l'égaliser par une flèche $j_\alpha\longrightarrow\bar{j}_\alpha$ dans $J'$. Donc quitte à remplacer $j_\alpha$ par % [sic] \varphi(i_\alpha) -> \varphi(i_{\alpha_0}), both composites of the double arrow start at \varphi(i_{\alpha_0})
\msp{page 324}
$\bar{j}_\alpha$, $u_\alpha$ par le composé $\bar{u}_\alpha:\varphi(i_0)\longrightarrow j_\alpha\longrightarrow\bar{j}_\alpha$ et $\tau'_{\alpha,\alpha_0}$ par le composé $\bar{\tau}'_{\alpha,\alpha_0}:j_{\alpha_0}\longrightarrow j_\alpha\longrightarrow\bar{j}_\alpha$ (sans toucher à $i_\alpha$, $v_{\alpha_0}$, $\tau_{\alpha_0,\alpha}$), on peut obtenir la commutativité, d'où $3_\alpha(\beta)$. % [sic] \varphi(i_0) -> \varphi(i_\alpha), u_\alpha has source \varphi(i_\alpha); [sic] \tau_{\alpha_0,\alpha} -> \tau_{\alpha,\alpha_0}, the arrow is written \tau_{\alpha,\alpha_0} everywhere else

Reste la contrainte $4_\alpha(\beta)$. Pour ceci, considérons le diagramme composite
\[
\begin{tikzcd}
j_\beta \arrow[r, "v_\beta"] \arrow[d, "\tau'_{\alpha_0,\beta}"'] & \varphi(i_{\beta+1}) \arrow[d, "\varphi(\tau_{\alpha,\beta+1})"] \\
j_{\alpha_0} \arrow[r, "v_{\alpha_0}"] & \varphi(i_{\alpha_0+1})=\varphi(i_\alpha)\ .
\end{tikzcd}
\]

Si $\alpha_0$ n'est pas un ordinal limite, mais de la forme $\alpha_1+1$, on peut encore se ramener au cas où $\beta=\alpha_1$, en procédant comme ci-dessus. Mais on peut procéder dans tous les cas, sans un tel artifice (\,!), en considérant les double-flèches
\[
(*)\qquad\qquad
\begin{tikzcd}
j_\beta \arrow[r, shift left, "\varphi(\tau_{\alpha,\beta+1})v_\beta"] \arrow[r, shift right, "v_{\alpha_0}\tau'_{\alpha_0,\beta}"'] & \varphi(i_\alpha)
\end{tikzcd}
\]
% foot: Version 17 décembre 2025
\printedpage{123}
% head: A. Grothendieck : Dérivateurs, Ch. XVIII, Édition des Universités de Montpellier II et Paris VII
déduites du carré précédent. Comme
\msp{page 325}
l'ensemble de ces double-flèches est de cardinal $\card\alpha<\pi$, et que $J$ est grande devant $\pi$ (il suffirait qu'il le soit pour les cardinaux $\pi'<\pi$), on peut trouver une flèche $\varphi(i_\alpha)\xrightarrow{u}j$ dans $J$ qui égalise toutes ces double-flèches. Comme $\varphi$ est cofinal, on peut supposer que $j$ est de la forme $\varphi(i)$. Comme $\varphi$ est un foncteur cofinal de catégories filtrantes, il existe alors $i\xrightarrow{\rho}i'$ dans $I$, tel que le composé $\varphi(i_\alpha)\xrightarrow{u}\varphi(i)\xrightarrow{\varphi(\rho)}\varphi(i')$ soit de la forme $\varphi(\sigma)$, pour $\sigma:i_\alpha\longrightarrow i'$. Au total, on a trouvé $\sigma:i_\alpha\longrightarrow i'$, tel que $\varphi(\sigma)$ $(=\varphi(\rho)u):\varphi(i_\alpha)\longrightarrow\varphi(i')$ égalise toutes les double-flèches envisagées dans $(*)$, p. 324. Donc il s'impose de remplacer $i_\alpha$ par $\bar{i}_\alpha=i'$, $v_{\alpha_0}$ par le composé $j_{\alpha_0}\longrightarrow\varphi(i_\alpha)\xrightarrow{\varphi(\sigma)}\varphi(\bar{i}_\alpha)$, $\tau_{\alpha,\alpha_0}$
\msp{page 326}
par le composé $\bar{\tau}_{\alpha,\alpha_0}:i_{\alpha_0}\xrightarrow{\tau_{\alpha,\alpha_0}}i_\alpha\xrightarrow{\sigma}\bar{i}_\alpha$.

Pour résumer, voici la recette pour construire, d'abord \ovalbox{$i_\alpha$, $v_{\alpha_0}$, $\tau_{\alpha,\alpha_0}$,}et \textit{ensuite} \ovalbox{$j_\alpha$, $u_\alpha$, $\tau'_{\alpha,\alpha_0}$}.

On procède d'abord comme page 321, 322 pour construire le premier triple, de façon à satisfaire à la première compatibilité $(*)$ (p. 321), i.e. à $1_\alpha$. On “corrige” ensuite ce choix de $i_\alpha$, $v_{\alpha_0}$, $\tau_{\alpha,\alpha_0}$, en remplaçant $i_\alpha$ par $\bar{i}_\alpha$ pour $\sigma:i_\alpha\longrightarrow\bar{i}_\alpha$, $v_{\alpha_0}$ par $\bar{v}_{\alpha_0}=\varphi(\sigma)v_{\alpha_0}$, et $\tau_{\alpha,\alpha_0}$ par $\bar{\tau}_{\alpha,\alpha_0}=\sigma\tau_{\alpha,\alpha_0}$, de sorte à obtenir les compatibilités $4_\alpha(\beta)$ pour $\beta<\alpha_0$, comme on vient de l'indiquer. On ne touchera plus à ces choix de $i_\alpha$, $v_{\alpha_0}$, $\tau_{\alpha,\alpha_0}$, dans la suite de la construction. On construit $(j_\alpha,u_\alpha,\tau'_{\alpha,\alpha_0})$ de façon à satisfaire à la deuxième compatibilité $(*)$ (p. 321), comme
\msp{page 327}
indiqué page 322. Ensuite, on corrige ce choix, en remplaçant $j_\alpha$ par $\bar{j}_\alpha$, avec $\rho:j_\alpha\longrightarrow\bar{j}_\alpha$, et $u_\alpha$ par $\bar{u}_\alpha=\rho u_\alpha$, $\tau'_{\alpha,\alpha_0}$ par $\bar{\tau}'_{\alpha,\alpha_0}=\rho\tau'_{\alpha,\alpha_0}$, de façon à satisfaire aux contraintes $3_\alpha(\beta)$ de la page 318.

Cela achève la construction dans le cas où $\alpha=\alpha_0+1$.

\textit{Cas 2 B : $\alpha$ est ordinal limite.}

Il faut donc construire les données 1) 2) 3) de la page 317 (à l'exclusion de la donnée 4), de façon à satisfaire à la seule contrainte $3_\alpha(\beta)$.

On commence à construire $i_\alpha$ et les $\tau_{\beta,\alpha}$ dans $I$. C'est possible, car $I_\alpha=\{\beta\mid\beta<\alpha\}$ est de cardinal $\pi_\alpha<\pi$ – il suffit donc de supposer $I$ grande devant les cardinaux $\pi'<\pi$.

On construit ensuite de même $J_\alpha$ dans $J'$ et les $\tau'_{\beta,\alpha}$ $(\beta<\alpha)$, en utilisant le % [sic] J_\alpha -> j_\alpha, the object of J' is meant
\msp{page 328}
fait que $J'$ est grande devant $\card\alpha=\card I_\alpha$. On doit de plus construire une flèche
\[
u_\alpha\ :\ \varphi(i_\alpha)\ \longrightarrow\ j_\alpha\ .
\]
C'est possible, car $J'$ étant cofinale dans $J$, il existe $j'$ dans $J'$ et $\varphi(i_\alpha)\longrightarrow J'$, et $J'$ étant filtrante, il existe $\bar{j}_\alpha$ et des flèches % [sic] \varphi(i_\alpha)\longrightarrow J' -> \varphi(i_\alpha)\longrightarrow j', the arrow goes to the object j'
$\begin{tikzcd}[ampersand replacement=\&, baseline=(\tikzcdmatrixname-1-1.base)]
j_\alpha \arrow[rr] \& \& j_\alpha \\
\& j' \arrow[ur] \&
\end{tikzcd}$, donc quitte à remplacer $j_\alpha$ par $\bar{j}_\alpha$, les $\tau'_{\alpha,\beta}$ par des $\bar{\tau}'_{\alpha,\beta}$, on trouve bien une flèche $u_\alpha:\varphi(i_\alpha)\longrightarrow j_\alpha$. Il faut ensuite s'occuper des % [sic] target j_\alpha in the small diagram -> \bar{j}_\alpha, the two arrows go to the object \bar{j}_\alpha just introduced
% foot: Version 17 décembre 2025

\end{document}
